% Options for packages loaded elsewhere
% Options for packages loaded elsewhere
\PassOptionsToPackage{unicode}{hyperref}
\PassOptionsToPackage{hyphens}{url}
\PassOptionsToPackage{dvipsnames,svgnames,x11names}{xcolor}
%
\documentclass[
  13pt,
]{article}
\usepackage{xcolor}
\usepackage[margin=2cm]{geometry}
\usepackage{amsmath,amssymb}
\setcounter{secnumdepth}{-\maxdimen} % remove section numbering
\usepackage{iftex}
\ifPDFTeX
  \usepackage[T1]{fontenc}
  \usepackage[utf8]{inputenc}
  \usepackage{textcomp} % provide euro and other symbols
\else % if luatex or xetex
  \usepackage{unicode-math} % this also loads fontspec
  \defaultfontfeatures{Scale=MatchLowercase}
  \defaultfontfeatures[\rmfamily]{Ligatures=TeX,Scale=1}
\fi
\usepackage{lmodern}
\ifPDFTeX\else
  % xetex/luatex font selection
  \setmainfont[]{Times New Roman}
\fi
% Use upquote if available, for straight quotes in verbatim environments
\IfFileExists{upquote.sty}{\usepackage{upquote}}{}
\IfFileExists{microtype.sty}{% use microtype if available
  \usepackage[]{microtype}
  \UseMicrotypeSet[protrusion]{basicmath} % disable protrusion for tt fonts
}{}
\usepackage{setspace}
\makeatletter
\@ifundefined{KOMAClassName}{% if non-KOMA class
  \IfFileExists{parskip.sty}{%
    \usepackage{parskip}
  }{% else
    \setlength{\parindent}{0pt}
    \setlength{\parskip}{6pt plus 2pt minus 1pt}}
}{% if KOMA class
  \KOMAoptions{parskip=half}}
\makeatother
% Make \paragraph and \subparagraph free-standing
\makeatletter
\ifx\paragraph\undefined\else
  \let\oldparagraph\paragraph
  \renewcommand{\paragraph}{
    \@ifstar
      \xxxParagraphStar
      \xxxParagraphNoStar
  }
  \newcommand{\xxxParagraphStar}[1]{\oldparagraph*{#1}\mbox{}}
  \newcommand{\xxxParagraphNoStar}[1]{\oldparagraph{#1}\mbox{}}
\fi
\ifx\subparagraph\undefined\else
  \let\oldsubparagraph\subparagraph
  \renewcommand{\subparagraph}{
    \@ifstar
      \xxxSubParagraphStar
      \xxxSubParagraphNoStar
  }
  \newcommand{\xxxSubParagraphStar}[1]{\oldsubparagraph*{#1}\mbox{}}
  \newcommand{\xxxSubParagraphNoStar}[1]{\oldsubparagraph{#1}\mbox{}}
\fi
\makeatother


\usepackage{longtable,booktabs,array}
\usepackage{calc} % for calculating minipage widths
% Correct order of tables after \paragraph or \subparagraph
\usepackage{etoolbox}
\makeatletter
\patchcmd\longtable{\par}{\if@noskipsec\mbox{}\fi\par}{}{}
\makeatother
% Allow footnotes in longtable head/foot
\IfFileExists{footnotehyper.sty}{\usepackage{footnotehyper}}{\usepackage{footnote}}
\makesavenoteenv{longtable}
\usepackage{graphicx}
\makeatletter
\newsavebox\pandoc@box
\newcommand*\pandocbounded[1]{% scales image to fit in text height/width
  \sbox\pandoc@box{#1}%
  \Gscale@div\@tempa{\textheight}{\dimexpr\ht\pandoc@box+\dp\pandoc@box\relax}%
  \Gscale@div\@tempb{\linewidth}{\wd\pandoc@box}%
  \ifdim\@tempb\p@<\@tempa\p@\let\@tempa\@tempb\fi% select the smaller of both
  \ifdim\@tempa\p@<\p@\scalebox{\@tempa}{\usebox\pandoc@box}%
  \else\usebox{\pandoc@box}%
  \fi%
}
% Set default figure placement to htbp
\def\fps@figure{htbp}
\makeatother


% definitions for citeproc citations
\NewDocumentCommand\citeproctext{}{}
\NewDocumentCommand\citeproc{mm}{%
  \begingroup\def\citeproctext{#2}\cite{#1}\endgroup}
\makeatletter
 % allow citations to break across lines
 \let\@cite@ofmt\@firstofone
 % avoid brackets around text for \cite:
 \def\@biblabel#1{}
 \def\@cite#1#2{{#1\if@tempswa , #2\fi}}
\makeatother
\newlength{\cslhangindent}
\setlength{\cslhangindent}{1.5em}
\newlength{\csllabelwidth}
\setlength{\csllabelwidth}{3em}
\newenvironment{CSLReferences}[2] % #1 hanging-indent, #2 entry-spacing
 {\begin{list}{}{%
  \setlength{\itemindent}{0pt}
  \setlength{\leftmargin}{0pt}
  \setlength{\parsep}{0pt}
  % turn on hanging indent if param 1 is 1
  \ifodd #1
   \setlength{\leftmargin}{\cslhangindent}
   \setlength{\itemindent}{-1\cslhangindent}
  \fi
  % set entry spacing
  \setlength{\itemsep}{#2\baselineskip}}}
 {\end{list}}
\usepackage{calc}
\newcommand{\CSLBlock}[1]{\hfill\break\parbox[t]{\linewidth}{\strut\ignorespaces#1\strut}}
\newcommand{\CSLLeftMargin}[1]{\parbox[t]{\csllabelwidth}{\strut#1\strut}}
\newcommand{\CSLRightInline}[1]{\parbox[t]{\linewidth - \csllabelwidth}{\strut#1\strut}}
\newcommand{\CSLIndent}[1]{\hspace{\cslhangindent}#1}



\setlength{\emergencystretch}{3em} % prevent overfull lines

\providecommand{\tightlist}{%
  \setlength{\itemsep}{0pt}\setlength{\parskip}{0pt}}



 


\usepackage[noblocks]{authblk}
\renewcommand*{\Authsep}{, }
\renewcommand*{\Authand}{, }
\renewcommand*{\Authands}{, }
\renewcommand\Affilfont{\small}
\makeatletter
\@ifpackageloaded{caption}{}{\usepackage{caption}}
\AtBeginDocument{%
\ifdefined\contentsname
  \renewcommand*\contentsname{Table of contents}
\else
  \newcommand\contentsname{Table of contents}
\fi
\ifdefined\listfigurename
  \renewcommand*\listfigurename{List of Figures}
\else
  \newcommand\listfigurename{List of Figures}
\fi
\ifdefined\listtablename
  \renewcommand*\listtablename{List of Tables}
\else
  \newcommand\listtablename{List of Tables}
\fi
\ifdefined\figurename
  \renewcommand*\figurename{Figure}
\else
  \newcommand\figurename{Figure}
\fi
\ifdefined\tablename
  \renewcommand*\tablename{Table}
\else
  \newcommand\tablename{Table}
\fi
}
\@ifpackageloaded{float}{}{\usepackage{float}}
\floatstyle{ruled}
\@ifundefined{c@chapter}{\newfloat{codelisting}{h}{lop}}{\newfloat{codelisting}{h}{lop}[chapter]}
\floatname{codelisting}{Listing}
\newcommand*\listoflistings{\listof{codelisting}{List of Listings}}
\makeatother
\makeatletter
\makeatother
\makeatletter
\@ifpackageloaded{caption}{}{\usepackage{caption}}
\@ifpackageloaded{subcaption}{}{\usepackage{subcaption}}
\makeatother
\usepackage{bookmark}
\IfFileExists{xurl.sty}{\usepackage{xurl}}{} % add URL line breaks if available
\urlstyle{same}
\hypersetup{
  pdftitle={Preferences for Market-Based Welfare: Disentangling the role of merit and privilege beliefs},
  pdfauthor={Andreas Laffert; Juan Carlos Castillo; René Canales; Tomás Urzúa},
  colorlinks=true,
  linkcolor={blue},
  filecolor={Maroon},
  citecolor={Blue},
  urlcolor={Blue},
  pdfcreator={LaTeX via pandoc}}


\title{Preferences for Market-Based Welfare: Disentangling the role of
merit and privilege beliefs}


\author{Andreas Laffert, Juan Carlos Castillo, René Canales, Tomás
Urzúa}
\affil{Department of Sociology, University of Chile}

\date{May 22, 2026}

\usepackage{xcolor}
\definecolor{mylinkcolor}{HTML}{0563C1}
\hypersetup{
  colorlinks = true,
  linkcolor  = mylinkcolor,
  citecolor  = mylinkcolor,
  urlcolor   = mylinkcolor
}
\begin{document}
\maketitle


\setstretch{1.5}
Word count: 451

\textbf{Background and relevance}. The commodification of social welfare
has reshaped not only how services are delivered but also the moral
frameworks through which citizens evaluate who deserves what and why
(\citeproc{ref-koos_moral_2019}{Koos \& Sachweh, 2019};
\citeproc{ref-mau_inequality_2015}{Mau, 2015};
\citeproc{ref-vanoorschot_social_2017}{Van Oorschot et al., 2017}).
Comparative welfare state research has increasingly documented that
marketization generates policy feedback effects on public attitudes,
strengthening support for market-based allocation of core services such
as health, education, and pensions; a phenomenon captured by the concept
of market justice preferences
(\citeproc{ref-busemeyer_welfare_2020}{Busemeyer \& Iversen, 2020};
\citeproc{ref-castillo_justification_2026}{Castillo et al., 2026},
\citeproc{ref-castillo_perceptions_2025}{2025};
\citeproc{ref-immergut_it_2020}{Immergut \& Schneider, 2020};
\citeproc{ref-lindh_public_2015}{Lindh, 2015};
\citeproc{ref-svallfors_political_2007}{Svallfors, 2007};
\citeproc{ref-zhu_policy_2015}{Zhu \& Lipsmeyer, 2015}). Within this
moral economy framework (\citeproc{ref-svallfors_moral_2006}{Svallfors,
2006}), criteria of deservingness play a central role in mediating the
ways in which structural conditions shape welfare attitudes. Drawing on
the NICER framework (\citeproc{ref-knotz_recast_2022}{Knotz et al.,
2022}), this paper focuses on the effort dimension of deservingness,
operationalized through meritocratic beliefs, as a key normative driver
of preferences for market justice. Chile offers a paradigmatic case:
decades of deep welfare commodification have institutionalized
individual responsibility as the dominant principle in pensions,
education, and health (\citeproc{ref-madariaga_three_2020}{Madariaga,
2020}; \citeproc{ref-valenzuela_socioeconomic_2013}{Valenzuela et al.,
2013}), making it an especially fertile context for examining how
meritocracy legitimizes market-based access to welfare.

\textbf{Research question and objective}. This study examines how
multidimensional meritocratic beliefs, encompassing both meritocratic
and non-meritocratic (privilege-based) perceptions and preferences,
shape support for market-based welfare across three domains: health,
education, and pensions. A key contribution is the explicit inclusion of
non-meritocratic beliefs (i.e., perceptions and endorsements of
privilege as a mechanism of social advancement), which the existing
literature has largely sidelined.

\textbf{Data and methods}. We analyze cross-sectional data from a survey
of urban Chilean adults (\(N\) = 2,557). Market justice preferences are
measured using items that assess the perceived fairness of income-based
access to health, education, and pensions. Meritocratic beliefs are
operationalized following a multidimensional approach
(\citeproc{ref-castillo_multidimensional_2023}{Castillo et al., 2023}),
distinguishing between meritocratic (effort and talent are rewarded) and
non-meritocratic (privilege and connections drive success) perceptions
and preferences. Structural equation modeling (SEM) is used to estimate
the associations between these four belief dimensions and market justice
preferences (\citeproc{ref-kline_principles_2023}{Kline, 2023}).

\textbf{Findings and implications}. Results show that meritocratic
perceptions are positively associated with market justice preferences,
replicating prior findings
(\citeproc{ref-castillo_justification_2026}{Castillo et al., 2026}).
Non-meritocratic perceptions (perceiving privilege) are negatively
associated, suggesting that awareness of unearned advantage undermines
support for market allocation. Meritocratic preferences show no
significant effect, possibly due to limited variance. Notably, the
strongest effect is observed for non-meritocratic preferences:
individuals who endorse privilege as a legitimate basis for social
reward show significantly higher market-justice preferences, the largest
effect in the model. These findings extend our understanding of how
deservingness beliefs operate across welfare domains, and highlight that
embracing non-meritocratic principles does not necessarily imply
opposition to market-based welfare. Methodologically, the study advances
the importance of disaggregating meritocratic beliefs into normative and
descriptive dimensions in the study of welfare attitudes.

\subsection{References}\label{references}

\phantomsection\label{refs}
\begin{CSLReferences}{1}{0}
\bibitem[\citeproctext]{ref-busemeyer_welfare_2020}
Busemeyer, M., \& Iversen, T. (2020). The {Welfare State} with {Private
Alternatives}: {The Transformation} of {Popular Support} for {Social
Insurance}. \emph{The Journal of Politics}, \emph{82}(2), 671--686.
\url{https://doi.org/10.1086/706980}

\bibitem[\citeproctext]{ref-castillo_justification_2026}
Castillo, J. C., Canales-Sellés, R., Laffert, A., \& Urzúa, T. (2026).
Justification trajectories for pension inequality in {Chile}
(2016--2023): The role of social class and beliefs in meritocracy.
\emph{Frontiers in Sociology}, \emph{11}, 1771856.
\url{https://doi.org/10.3389/fsoc.2026.1771856}

\bibitem[\citeproctext]{ref-castillo_multidimensional_2023}
Castillo, J. C., Iturra, J., Maldonado, L., Atria, J., \& Meneses, F.
(2023). A {Multidimensional Approach} for {Measuring Meritocratic
Beliefs}: {Advantages}, {Limitations} and {Alternatives} to the {ISSP
Social Inequality Survey}. \emph{International Journal of Sociology},
\emph{53}(6), 448--472.
\url{https://doi.org/10.1080/00207659.2023.2274712}

\bibitem[\citeproctext]{ref-castillo_perceptions_2025}
Castillo, J. C., Laffert, A., Carrasco, K., \& Iturra-Sanhueza, J.
(2025). Perceptions of inequality and meritocracy: Their interplay in
shaping preferences for market justice in {Chile} (2016--2023).
\emph{Frontiers in Sociology}, \emph{10}, 1634219.
\url{https://doi.org/10.3389/fsoc.2025.1634219}

\bibitem[\citeproctext]{ref-immergut_it_2020}
Immergut, E. M., \& Schneider, S. M. (2020). Is it unfair for the
affluent to be able to purchase {``better''} healthcare? {Existential}
standards and institutional norms in healthcare attitudes across 28
countries. \emph{Social Science \& Medicine}, \emph{267}, 113146.
\url{https://doi.org/10.1016/j.socscimed.2020.113146}

\bibitem[\citeproctext]{ref-kline_principles_2023}
Kline, R. B. (2023). \emph{Principles and {Practice} of {Structural
Equation Modeling}}. Guilford Publications.

\bibitem[\citeproctext]{ref-knotz_recast_2022}
Knotz, C. M., Gandenberger, M. K., Fossati, F., \& Bonoli, G. (2022). A
{Recast Framework} for {Welfare Deservingness Perceptions}. \emph{Social
Indicators Research}, \emph{159}(3), 927--943.
\url{https://doi.org/10.1007/s11205-021-02774-9}

\bibitem[\citeproctext]{ref-koos_moral_2019}
Koos, S., \& Sachweh, P. (2019). The moral economies of market
societies: Popular attitudes towards market competition, redistribution
and reciprocity in comparative perspective. \emph{Socio-Economic
Review}, \emph{17}(4), 793--821.
\url{https://doi.org/10.1093/ser/mwx045}

\bibitem[\citeproctext]{ref-lindh_public_2015}
Lindh, A. (2015). Public {Opinion} against {Markets}? {Attitudes}
towards {Market Distribution} of {Social Services} -- {A Comparison} of
17 {Countries}. \emph{Social Policy \& Administration}, \emph{49}(7),
887--910. \url{https://doi.org/10.1111/spol.12105}

\bibitem[\citeproctext]{ref-madariaga_three_2020}
Madariaga, A. (2020). The three pillars of neoliberalism: {Chile}'s
economic policy trajectory in comparative perspective.
\emph{Contemporary Politics}, \emph{26}(3), 308--329.
\url{https://doi.org/10.1080/13569775.2020.1735021}

\bibitem[\citeproctext]{ref-mau_inequality_2015}
Mau, S. (2015). \emph{Inequality, {Marketization} and the {Majority
Class}: {Why Did} the {European Middle Classes Accept Neo-Liberalism}?}
London: Palgrave Macmillan UK.

\bibitem[\citeproctext]{ref-svallfors_moral_2006}
Svallfors, S. (2006). \emph{The moral economy of class}. Stanford
University Press.

\bibitem[\citeproctext]{ref-svallfors_political_2007}
Svallfors, S. (Ed.). (2007). \emph{The {Political Sociology} of the
{Welfare State}: {Institutions}, {Social Cleavages}, and {Orientations}}
(1st ed.). Stanford University Press.
\url{https://doi.org/10.2307/j.ctvr0qv0q}

\bibitem[\citeproctext]{ref-valenzuela_socioeconomic_2013}
Valenzuela, J. P., Bellei, C., \& Ríos, D. D. L. (2013). Socioeconomic
school segregation in a market-oriented educational system. {The} case
of {Chile}. \emph{Journal of Education Policy}, \emph{29}(2), 217--241.
\url{https://doi.org/10.1080/02680939.2013.806995}

\bibitem[\citeproctext]{ref-vanoorschot_social_2017}
Van Oorschot, W., Roosma, F., Meuleman, B., \& Reeskens, T. (Eds.).
(2017). \emph{The {Social Legitimacy} of {Targeted Welfare}: {Attitudes}
to {Welfare Deservingness}}. Edward Elgar Publishing.
\url{https://doi.org/10.4337/9781785367212}

\bibitem[\citeproctext]{ref-zhu_policy_2015}
Zhu, L., \& Lipsmeyer, C. S. (2015). Policy feedback and economic risk:
The influence of privatization on social policy preferences.
\emph{Journal of European Public Policy}, \emph{22}(10), 1489--1511.
\url{https://doi.org/10.1080/13501763.2015.1031159}

\end{CSLReferences}




\end{document}
